\chapter{Gaussian Mixture Models evaluation}
\label{chap:gmm_evaluation}

In this chapter, we evaluate the performance of Gaussian Mixture Models (GMMs) on the project dataset. Unlike the standard Multivariate Gaussian (MVG) classifier which assumes each class can be modeled by a single normal distribution, GMMs allow us to approximate any sufficiently regular unknown density by combining multiple Gaussian components.

Following the theoretical framework, we independently fit one GMM to describe the density of the "Genuine" class and one GMM to describe the density of the "Impostor" class. The class-conditional density for each class is modeled as a weighted sum of $M$ Gaussians:
\begin{equation}
    f_X(x) = \sum_{g=1}^M w_g \mathcal{N}(x | \mu_g, \Sigma_g)
\end{equation}
. At inference time, we compute the class posterior probabilities using these GMMs as the class-conditional densities.

\section{Experimental Results}

We trained GMMs for both classes varying the number of components ($M \in \{1, 2, 4, 8, 16\}$) across three covariance architectures: Full, Diagonal, and Tied. The models were evaluated on the evaluation set for the primary application ($\pi_T = 0.1$). 

The Minimum DCF and Actual DCF metrics are summarized in Table \ref{tab:gmm_results}.

\begin{table}[H]
    \centering
    \caption{GMM performance (minDCF / actDCF) for $\pi_T = 0.1$}
    \label{tab:gmm_results}
    \resizebox{\textwidth}{!}{
    \begin{tabular}{lccccc}
        \toprule
        \textbf{Covariance type} & \textbf{1 Comp.} & \textbf{2 Comp.} & \textbf{4 Comp.} & \textbf{8 Comp.} & \textbf{16 Comp.} \\
        \midrule
        \textbf{Full Covariance} & 0.1302 / 0.1399 & 0.1342 / 0.1400 & 0.0720 / 0.0730 & 0.0659 / 0.0660 & 0.0620 / 0.0640 \\
        \textbf{Diagonal Covariance} & 0.1311 / 0.1439 & 0.1360 / 0.1479 & 0.0610 / 0.0661 & 0.0610 / 0.0681 & \textbf{0.0579} / \textbf{0.0599} \\
        \textbf{Tied Covariance} & 0.1302 / 0.1399 & 0.1302 / 0.1399 & 0.1302 / 0.1399 & 0.1342 / 0.1449 & 0.0900 / 0.0910 \\
        \bottomrule
    \end{tabular}
    }
\end{table}

\section{Analysis and theoretical considerations}

\subsection*{Question: What do you observe as the number of components increases?}

For the Full and Diagonal covariance models, there is a massive leap in discriminatory power as the number of components increases from 2 to 4. For instance, the Full Covariance minDCF drops precipitously from $0.1342$ (2 components) to $0.0720$ (4 components), and continues to improve up to $0.0620$ at 16 components. 

\textbf{Theoretical explanation:} This observation perfectly aligns with the fundamental property of GMMs. Because a single Gaussian (1 component) cannot adequately capture the heavy skewness of Features 5 and 6 (which we discovered in exploratory analysis), its performance stalls. By increasing the number of components $M$, the GMM successfully approximates the true, complex, non-Gaussian probability density of the data. The sharp drop in error at $M=4$ suggests that the underlying feature distributions contain significant multi-modality or severe asymmetries that require at least 4 local Gaussian components to be mapped accurately.

\subsection*{Question: Which covariance architecture works best and why?}

The \textbf{Diagonal Covariance GMM with 16 components} achieves the absolute best performance across all tested configurations, yielding a minDCF of \textbf{0.0579} and an exceptionally well-calibrated actDCF of \textbf{0.0599}. 

Conversely, the Tied Covariance model performs the worst, stalling at a minDCF of $0.0900$ even with 16 components. 

\textbf{Theoretical explanation:} We can explain this hierarchy by looking at the interaction between the mathematical properties of the models and the physical characteristics of our dataset:

\begin{itemize}
    \item \textbf{Why Diagonal wins over Full:} While a Full covariance matrix captures correlations between features, it requires estimating a massive number of parameters ($D(D+1)/2$ per component). As stated in the theoretical lectures, diagonal covariance matrices can approximate any distribution provided that enough components are used. Furthermore, in Laboratory 05, we mathematically proved that the features in this specific dataset possess near-zero correlation. Because the features are already inherently independent, computing the full cross-covariance terms only introduces parameter estimation noise. The Diagonal GMM successfully avoids this parameter explosion, serving as an effective regularizer while using its 16 independent components to perfectly map the density contours.
    
    \item \textbf{Why Tied Covariance fails:} The Tied GMM forces all components within a class (or across classes, depending on the implementation) to share the exact same covariance matrix, varying only their means. We know from our dataset exploration that the "Genuine" and "Impostor" classes have fundamentally different geometric spreads (e.g., Feature 1 variance mismatch). By restricting the components to a shared covariance geometry, the model loses the structural flexibility required to accurately bound the complex shape of the target class, resulting in the highest error rates among the three variants.
\end{itemize}

\subsection*{Question: Are the models well-calibrated?}

Yes, unlike the highly-regularized discriminative models from previous laboratories (such as SVMs or Quadratic LR), the GMMs exhibit spectacular native calibration. 

For the optimal Diagonal model at 16 components, the gap between the Minimum DCF ($0.0579$) and the Actual DCF ($0.0599$) is a negligible $0.0020$. Because the GMM is a generative model that explicitly estimates the true density $f_X(x)$, its raw output scores function as mathematically rigorous Log-Likelihood Ratios. Consequently, comparing these scores directly against the theoretical Bayes threshold for $\pi_T = 0.1$ yields near-optimal empirical decisions without requiring any post-hoc score calibration.